Write \(\mathcal O_a=\{O\in\mathcal O:O\subseteq\mathcal U_a\}\). For \(k_a\le M\),
\[
s_a(O_{T,a})\le Q_{k_a}(s_a(O_{1,a}),\ldots,s_a(O_{M,a}))\iff\sum_{i=1}^M\mathbf1\{s_a(O_{i,a})<s_a(O_{T,a})\}<k_a.
\tag{1}
\]
Thus ties are treated by strict comparison and the event depends only on the total preorder induced by \(s_a\). Conversely, represent ordered nonempty classes \(C_{a,1}\prec_a\cdots\prec_a C_{a,d_a}\) by score zero if \(d_a=1\), and otherwise by \((j-1)K/(d_a-1)\) on class \(j\). Every preorder has a representative in \([0,K]\) with at most \(m_a\) levels. The defining maximum is therefore attained by finite preorder enumeration, exactly over the declared score-map class.

If \(k_0,k_1\le M\), \(\texttt{lem:coherent-two-finite-rank-frontier-evaluator}\) gives for every preorder pair \(C\)
\[
U_C\le L_C\le\min\{L,d_0d_1\},\qquad \Pr\{\text{both rank events}\}=\sum_{\pi\in\Pi_C}\omega_C(\pi)D_{M,k_0,k_1}(\pi).
\tag{2}
\]
The exact key \((M,k_0,k_1,\pi_{00},\pi_{01},\pi_{10})\) determines the recurrence value, while \(\omega_C\) enters only the preorder objective. Local grouping, global balanced-tree lookup, table construction, and one computation for each of the \(G\) keys yield
\[
O\!\left(|\mathcal V|+F_{m_0}F_{m_1}(m_0m_1+L)+\sum_{C\in\mathfrak C}[L_C\log(1+L_C)+L_C+U_C]+\left\{\sum_CU_C\right\}\log(1+G)+GMk_0k_1\right).
\tag{3}
\]
Replacing \(L_C,U_C\) by their maxima gives the stated uniform bound; the tables use \(O(m_0m_1+L+k_0k_1+U_{\max}+G)\) memory. Every \(\pi\) comes from the unchanged \(R_w\), and every weight from the pushforward of the one \(P_T\), so memoization preserves both probability laws. The multinomial and corner-oriented lemmas compute the same per-key probability, and the rolling option proves the displayed coarse resource charges.

If exactly one rank \(k_b=M+1\), put \(c=1-b\). The boundary event is certain. Conditional on \(O_{T,c}=O_t\), (1) is a binomial count with
\[
\theta(O_t)=\sum_{O:s_c(O)<s_c(O_t)}q_{c,O},
\]
so target independence gives
\[
\sum_{O_t}p_{c,O_t}\Pr\{\operatorname{Bin}(M,\theta(O_t))<k_c\}.
\tag{4}
\]
For tie classes \(C_1\prec\cdots\prec C_d\), (4) is \(\sum_jp_j\Pr\{\operatorname{Bin}(M,\theta_j)<k_c\}\), where \(q_j=q_c(C_j)\), \(p_j=p_c(C_j)\), and \(\theta_j=\sum_{\ell<j}q_\ell\). A single class scan forms these quantities. For \(0<\theta<1\), consecutive binomial masses satisfy
\[
t_0=(1-\theta)^M,\quad t_{r+1}=t_r\frac{M-r}{r+1}\frac{\theta}{1-\theta},
\qquad\nu_M=\theta^M,\quad u_{r-1}=u_r\frac r{M-r+1}\frac{1-\theta}{\theta}.
\tag{5}
\]
Exponentiation by squaring initializes an endpoint in \(O(\log M)\), and the shorter tail uses \(O(\min\{k_c,M-k_c+1\})\) additional operations. At \(\theta=0,1\), the probabilities are one and zero without division. The one-boundary evaluator consequently gives the asserted \(F_{m_c}\)-preorder time and constant auxiliary CDF memory. If both ranks equal \(M+1\), both events are certain and the frontier equals \([(1-\alpha)-1]_+=0\).

For the nonbinary certificate, take four pair blocks, \(\mathcal G=S_2^4\), fair independent within-pair assignments, and a singleton selector. Let each arm's equivariant kernel choose a block with
\[
q=\left(\frac{13}{20},\frac1{20},\frac14,\frac1{20}\right)
\]
and select that block's unique observed arm member. Let the common target choose a block with
\[
p=\left(\frac1{20},\frac1{10},\frac14,\frac35\right)
\]
and split equally within it. Take \(M=3\), \(\alpha=1/2\), and \(\varepsilon_0=\varepsilon_1=1/4\), so both ranks equal three. Put \(s_0\equiv0\) and order arm-one blocks as \(1\prec3\prec2\prec4\). Its failure probability is
\[
\frac14\left(\frac{13}{20}\right)^3+\frac1{10}\left(\frac{18}{20}\right)^3+\frac35\left(\frac{19}{20}\right)^3=\frac{104957}{160000},
\tag{6}
\]
so the objective is \(24957/160000\). For a binary pair, let \(H_a\) be arm \(a\)'s high-score block set and \(A_a=(1-q(H_a))^3\). Inclusion--exclusion gives
\[
F_{\mathrm{bin}}(H_0,H_1)=p(H_0)A_0+p(H_1)A_1-p(H_0\cap H_1)A_0A_1.
\tag{7}
\]
In the row order \(\varnothing,1,2,3,4,12,13,14,23,24,34,123,124,134,234,1234\), exact substitution of all sixteen \(H_1\) choices gives row-maximum numerators
\[
\begin{split}
(&164616000,165302000,192052000,198366000,207906936,165912000,164712000,165787368,\\
&203032000,208582976,208582976,164632000,165793875,164631423,202894331,164616000)
\end{split}
\tag{8}
\]
over denominator \(320000000\). This exhausts all \(16^2\) binary pairs. Their largest objective is \(759109/5000000\), and the four-level advantage is
\[
\frac{24957}{160000}-\frac{759109}{5000000}=\frac{83189}{20000000}>0.
\tag{9}
\]
Finally, under disjoint target and calibration supports in an active arm, assign score \(K\) to its target-supported orbits and zero elsewhere, and score zero throughout the other arm. One rank event fails surely, so the deficit reaches its maximum \(1-\alpha\). At a boundary rank, that arm's event is certain and (4) is exactly the remaining arm's event. On the equal-split diagonal, \(k_0=k_1=k\), yielding \(\texttt{thm:sharp-orbit-coverage-frontier}\).